We present~BEAMFS, a recovery calculus for filesystem resilience under
Single-Event Effects (SEE). The work is the formal counterpart to the
FTRFS implementation experience reported in~\cite{desbrieres2026ftrfs}:
the empirical limits encountered during the FTRFS validation campaign
motivate a mathematical treatment of \emph{recovery itself} as a
first-class object, rather than as a side-effect of detection-and-correction
loops layered onto an unmodelled state space.

We define a recovery operator~\(\mathcal{G}\) over the joint state space
of a filesystem's volatile and persistent memory layers, formulate the
soundness property as a fixpoint convergence theorem under physically
realistic SEE perturbations, and derive an algorithm whose Linux-kernel
implementation path is sketched. Three memory technology layers are
treated explicitly: \mbox{DRAM} (with a forward note on
\mbox{DDR5} on-die~\mbox{ECC}), \mbox{SRAM}, and \mbox{NAND} flash via
\mbox{NVMe}. Emerging non-volatile technologies (\mbox{3D~XPoint},
\mbox{ReRAM}, \mbox{MRAM}) are positioned but left out of v1's modelling.

This Technical Report v1 establishes authorship of the approach,
fixes its scope, and lays the mathematical scaffolding for a reference
implementation in the planned successor to FTRFS. The empirical
section reinterprets the four bench metrics of the FTRFS v1
validation through the recovery-calculus lens, identifying which
metrics are actually \emph{recovery-bounded} and which are not.
No accelerator-beam measurements are claimed; the present document
is intended, among other purposes, as the scientific dossier from
which beam-time at facilities such as RADEF, GANIL or TRIUMF would
later be requested.
